\section{Resource description}
HCMO uses the persistent ontology IRI \nolinkurl{https://w3id.org/hcmo/ontology/hcm} and the stable term namespace \nolinkurl{https://w3id.org/hcmo/ontology/hcm#}. Domain-specific terms use four sub-namespaces beneath the same base: \nolinkurl{hcm/bio#}, \nolinkurl{hcm/env#}, \nolinkurl{hcm/obs#}, and \nolinkurl{hcm/tech#}. Published IRIs are not renamed when terminology changes; obsolete IRIs are retained in a compatibility module with deprecation and replacement metadata where a defensible replacement exists.

\textbf{Modular organisation.} The paper-matching HCMO candidate is authored as five active modules. The minimal core centres on \nolinkurl{hcm:Enclosure}, \nolinkurl{hcm:MonitoredEnclosure}, enclosure dimensions, enrichment, and stable housing relations. The bio module represents subjects, experimental groups, housing assignments, and study factors. The env module represents environmental profiles, light cycles, gas profiles, measurement specifications, and environmental properties. The obs module owns SOSA-style observations and their categorical, quantitative, behavioral, and tabular results. The tech module represents sensors, actuators, hardware, software, and time-series resources. A sixth, migration-only module preserves deprecated HCMO 0.0.1 IRIs but is not a source of terms for new data.

This organisation preserves two boundaries that are important for HCM data. First, housing context is not reduced to the animal: a housing assignment is an explicit record linking a subject or group to an enclosure. Second, a physical sensor, the observation it performs, and the result it produces are separate entities. For example, \nolinkurl{hcm-tech:Sensor} is linked to a monitored enclosure, captures a \nolinkurl{sosa:ObservableProperty}, and may make a \nolinkurl{sosa:Observation}; the observation then identifies its feature of interest and result using SOSA relations. This supports provenance-sensitive queries without treating a device, event, and data value as the same thing.

Environmental profiles, measurement specifications, and observations likewise use different relations: a profile composes specifications, each specification identifies the property and optional target quantity, and each observation uses the SOSA observed-property/result pattern. Housing membership is derived from time-bounded assignments at a stated time. Operational and calibration states are time-bounded records generated by explicit PROV activities rather than timeless booleans.

\textbf{Standards reuse.} Physical subjects, enclosures, enrichments, sensors, actuators, hardware, and experimental groups are exposed under BFO material entity in the default end-user hierarchy. Records, specifications, profiles, software, and observation-result artifacts use the IAO information-content anchor. A separate optional developer profile restores BFO's source-faithful intermediate hierarchy and the more precise object-aggregate classification of experimental groups. Sensors, actuators, observations, results, and observed properties reuse SOSA classes and relations, and \nolinkurl{hcm-tech:captures} is explicitly a subproperty of \nolinkurl{sosa:observes}. Schema.org supplies contributor and place exchange types. The example ABox uses OWL-Time intervals and the duration competency query consumes this representation. These are selective references, not claims that HCMO imports or reproduces the full external ontologies.

Numeric dimensions, sampling rates, environmental targets, and observation results use the pinned QUDT 3.4.0 \nolinkurl{QuantityValue}, \nolinkurl{numericValue}, and \nolinkurl{hasUnit} pattern \cite{qudt}. HCMO role properties state what each quantity describes. This is selective reuse of the checksummed schema and unit vocabulary, not a claim of full QUDT conformance.

\textbf{Claim strength.} We distinguish implemented semantic reuse or alignment, validated interoperability evidence, and formal profile conformance. Implemented alignment requires canonical versioned terms, reviewed semantic fit, and a normative role in HCMO; its assertion strength is stated per term and never generalized to a whole vocabulary. Interoperability evidence means that a pinned mixed-vocabulary example is validated and queried, but does not by itself establish ontology mappings or complete round trips. Formal conformance is reserved for a declared version and scope in which every applicable normative profile requirement has been tested. ``Provisional'' identifies work that has not yet qualified as implemented alignment; it is not a conformance level.

HCMO normatively follows the 2017 SOSA/SSN Recommendation \cite{sosa}, pinned to an immutable copy of the 2017 ontology artifact. Environmental and sensor-captured properties use \nolinkurl{sosa:ObservableProperty}; terms introduced by the developing later edition are not mixed into this release. HCMO also selectively reuses canonical unversioned SemTS 1.2.0 terms \cite{semts}: a location result table is a \nolinkurl{semts:TimeSeriesSegment} whose dimensions may be linked with \nolinkurl{semts:segmentDimension} to \nolinkurl{semts:DataDimension}. The earlier observation dimension restriction and knowledge-generation-specific \nolinkurl{semts:generated} relation were removed after semantic review. A pinned example-level evidence slice now uses PROV-O, specific OBI and STATO types, and ISA/Bioschemas exchange terms in one acyclic workflow \cite{isa,rocrate}. This is validated interoperability evidence, not an HCMO class mapping or a claim of formal ISA RO-Crate conformance. A second generated 2 $\times$ 2 fixture preserves animal identity, time-bounded housing, Source/Sample roles, repeated observations, factor/group structure, and separate STATO results and file fragments across a graph-isomorphic HCMO RDF/extended ISA RO-Crate pair. RO-Crate 1.2 and ISA-specific required validator rules pass. Formal ISA conformance remains deferred because the draft has no permanent profile URI and its validator currently disagrees about the base RO-Crate edition. ISA-JSON and ISA-Tab are controlled-loss projections. Their executable native overlap, tested with pinned \nolinkurl{isatools} 0.14.3, preserves a distinct animal Source, a real tissue Sample, their collection process and their HCMO identifiers through a JSON--Tab--JSON cycle. Housing, direct whole-animal recording, source-bound factor assignments, explicit groups, repeated observations, and semantic STATO/file-fragment links remain declared losses rather than being represented with fabricated Sample proxies. Broader logical mappings and integration remain separately reviewed work.

\textbf{Distribution and application support.} The authored modules are merged deterministically into Turtle, RDF/XML, and JSON-LD distributions. A generated \nolinkurl{profile.json} provides a flat term inventory for synchronization layers and interfaces, while \nolinkurl{ontology/context.jsonld} supports JSON-LD applications. \nolinkurl{hcmo.yaml} is the stable machine-readable contract that identifies the module order, generated artifacts, shapes, examples, and competency-query index.

The active source inventory contains 33 classes and 69 non-deprecated object or datatype properties. The generated release additionally contains deprecated compatibility terms and directly referenced external vocabulary terms. A signed class audit and property audit record the asserted semantics, restrictions, inference consequences, current evidence, and unresolved decisions without turning review notes into unattended ontology rewrites.

\begin{figure*}[t]
\centering
\resizebox{\textwidth}{!}{%
\begin{tikzpicture}[
  node distance=7mm and 9mm,
  entity/.style={draw, rounded corners=1.5pt, align=center, minimum height=9mm, font=\scriptsize, inner xsep=5pt},
  rel/.style={-{Latex[length=1.8mm]}, font=\tiny},
  evidence/.style={draw, dashed, rounded corners=1.5pt, align=center, font=\scriptsize, inner sep=4pt}
]
\node[entity, fill=yellow!14] (cage) {\texttt{hcm:MonitoredEnclosure}\\central physical hub};
\node[entity, fill=red!8, left=18mm of cage] (assignment) {\texttt{hcm-bio:HousingAssignment}\\time-bounded record};
\node[entity, fill=red!8, above=of assignment] (subject) {\texttt{hcm-bio:Subject}\\stable animal identity};
\node[entity, fill=violet!9, right=12mm of cage] (sensor) {\texttt{hcm-tech:Sensor}\\\texttt{sosa:Sensor}};
\node[entity, fill=blue!8, below=10mm of sensor] (observation) {HCMO observation\\\texttt{sosa:Observation}};
\node[entity, fill=blue!8, right=8mm of observation] (result) {\texttt{hcm-obs:}\\\texttt{ObservationResult} / \texttt{sosa:Result}};
\node[entity, fill=green!9, below=10mm of cage] (profile) {\texttt{hcm-env:}\\\texttt{EnvironmentProfile}};
\node[entity, fill=green!9, right=8mm of profile] (spec) {\texttt{hcm-env:}\\\texttt{MeasurementSpecification}};
\node[entity, fill=green!9, below=8mm of spec] (quantity) {\texttt{hcm-obs:QuantityValue}\\\texttt{qudt:QuantityValue}};
\node[evidence, below=of assignment] (time) {\texttt{time:Interval}\\assignment validity};

\draw[rel] (subject) -- node[left]{\texttt{hasHousingAssignment}} (assignment);
\draw[rel] (assignment) -- node[above]{\texttt{assignedToEnclosure}} (cage);
\draw[rel] (assignment) -- node[left]{\texttt{time:hasTime}} (time);
\draw[rel] (cage) -- node[above]{\texttt{monitoredBy}} (sensor);
\draw[rel] (sensor) -- node[right]{\texttt{sosa:madeObservation}} (observation);
\draw[rel] (observation) -- node[above]{\texttt{sosa:hasResult}} (result);
\draw[rel, dashed] (quantity) to[bend right=12] node[right]{subclass} (result);
\draw[rel] (observation) to[bend right=18] node[above]{\texttt{sosa:hasFeatureOfInterest}} (subject);
\draw[rel] (cage) -- node[left]{\texttt{hasEnvironment}} (profile);
\draw[rel] (profile) -- node[above]{\texttt{hasMeasurementSpec}} (spec);
\draw[rel] (spec) -- node[above]{\texttt{hasSpecifiedValue}} (quantity);
\end{tikzpicture}
}
\caption{The five-module HCMO model preserves animal identity, temporal housing,
environmental specifications, sensing events, and results as distinct entities.}
\label{fig:domain-model}
\end{figure*}

